\section{Investigación técnica (Defensa TFG): modelos adicionales, hiperparámetros y validación}

\subsection{1) Alcance}

Este documento complementa al informe de arquitectura y \textbf{no} desarrolla de nuevo:

\begin{itemize}
  \item esquema canónico (ya documentado)
  \item algoritmo QRDQN en detalle (ya documentado)
\end{itemize}

Aquí se resumen: otros modelos/algoritmos, parámetros y sistema de evaluación.

\subsection{2) Modelos y algoritmos presentes además del núcleo QRDQN}

\subsubsection{2.1 Baseline clásico: Random Forest}

Archivo: \verb|src/baseline_random_forest.py|

\begin{itemize}
  \item Modelo: \texttt{RandomForestClassifier}
  \item Hiperparámetros por defecto:
    \begin{itemize}
      \item \texttt{n\_estimators=200}
      \item \texttt{max\_depth=None}
      \item \texttt{n\_jobs=-1}
      \item \texttt{random\_state=42}
    \end{itemize}
  \item Se puede ejecutar sobre CICIDS2017 o NSL-KDD (histórico)
  \item Evalúa con \texttt{confusion\_matrix} y \texttt{classification\_report}
\end{itemize}

Uso en defensa: baseline no-RL para comparar la estrategia RL contra un método supervisado clásico.

\subsubsection{2.2 DQN como fallback operativo}

Archivo: \verb|scripts/predict_real_traffic_v2.py| (\texttt{load\_model}).

\begin{itemize}
  \item Intenta cargar con \texttt{sb3\_contrib.QRDQN}
  \item Si ocurre una excepción durante esa carga, hace fallback a \texttt{stable\_baselines3.DQN}
\end{itemize}

Esto da robustez operativa de inferencia cuando falla la carga de \texttt{QRDQN}; la ausencia de
\texttt{sb3\_contrib} es un caso posible, pero no el único.

\subsubsection{2.3 Búsqueda de hiperparámetros con Optuna}

Archivo: \verb|src/tune_hparams.py|

Espacio de búsqueda:

\begin{itemize}
  \item \texttt{learning\_rate}: \texttt{1e-5} \ldots \texttt{1e-3} (log)
  \item \texttt{batch\_size}: \texttt{\{256, 512, 1024, 2048\}}
  \item \texttt{gradient\_steps}: \texttt{\{10, 50, 100\}}
  \item \texttt{net\_arch}: \texttt{256\_128}, \texttt{512\_256}, \texttt{256\_256}
  \item \texttt{gamma}: \texttt{0.95} \ldots \texttt{0.999}
  \item \texttt{train\_freq}: \texttt{\{50, 100, 200\}}
\end{itemize}

Objetivo optimizado: \texttt{F1} de la clase ataque (\texttt{pos\_label=1}).

\subsection{3) Arquitecturas de red usadas en el proyecto}

Aunque la investigación específica de QRDQN está en otro documento, en el código aparecen
arquitecturas MLP concretas:

\begin{itemize}
  \item \texttt{net\_arch=[512,~256]}
    \begin{itemize}
      \item entrenamiento principal (\verb|src/train_rl_defender.py|)
      \item validación Check~C y leave-one-CSV-out
    \end{itemize}
  \item \texttt{net\_arch=[256,~128]}
    \begin{itemize}
      \item validación Check~B (labels barajadas)
    \end{itemize}
\end{itemize}

Interpretación para defensa:

\begin{itemize}
  \item la arquitectura más grande (\texttt{512,256}) se reserva para entrenamiento principal y
        validación exigente
  \item la más pequeña (\texttt{256,128}) acelera el test anti-leakage
\end{itemize}

\subsection{4) Hiperparámetros clave de entrenamiento y validación}

\subsubsection{4.1 Entrenamiento principal (\texttt{src/train\_rl\_defender.py})}

Comunes:

\begin{itemize}
  \item \texttt{learning\_rate=1e-4}
  \item \texttt{gamma=0.99}
  \item \texttt{tau=1.0}
  \item \texttt{buffer\_size=min(200\_000, max(total\_timesteps, 10\_000))}
\end{itemize}

Dependientes del preset:

\begin{itemize}
  \item \textbf{fast}
    \begin{itemize}
      \item \texttt{batch\_size=512}
      \item \texttt{gradient\_steps=10}
      \item \texttt{train\_freq=50}
      \item \texttt{target\_update\_interval=1\_000}
      \item \texttt{total\_timesteps} default: \texttt{25\_000}
    \end{itemize}
  \item \textbf{full}
    \begin{itemize}
      \item \texttt{batch\_size=2048}
      \item \texttt{gradient\_steps=20}
      \item \texttt{train\_freq=100}
      \item \texttt{target\_update\_interval=10\_000}
      \item \texttt{total\_timesteps} default: \texttt{100\_000}
    \end{itemize}
\end{itemize}

Split de datos:

\begin{itemize}
  \item \texttt{split\_mode=random} o \texttt{split\_mode=day}
  \item presets de filas (\texttt{fast} vs \texttt{full}) vía \verb|load_cicids2017_split|
\end{itemize}

\subsubsection{4.2 Leave-one-exact-CSV-out (\texttt{src/validate\_leave\_one\_csv\_out.py})}

\begin{itemize}
  \item \texttt{timesteps} por fold: default \texttt{30\_000}
  \item \texttt{predict\_batch\_size} default \texttt{8192}
  \item entrenamiento por fold con \texttt{net\_arch=[512,256]}, \texttt{batch\_size=512},
        \texttt{gradient\_steps=20}, \texttt{train\_freq=100}
  \item agrega métricas por fold y globales
\end{itemize}

\subsection{5) Métricas y validación experimental}

\subsubsection{5.1 Check A (evaluación directa)}

Archivo: \verb|src/validate_checks.py|, función \verb|check_a_direct_eval|

\begin{itemize}
  \item Predicción directa \verb|model.predict(X_test[i])|
  \item Reporta: accuracy, precision/recall/F1 para ataque y benigno, matriz de confusión y
        TP/FP/FN/TN
\end{itemize}

Objetivo: validar rendimiento sin depender de señales auxiliares del entorno.

\subsubsection{5.2 Check B (anti-leakage con etiquetas barajadas)}

Archivo: \verb|src/validate_checks.py|, función \verb|check_b_shuffled_labels|

\begin{itemize}
  \item baraja \texttt{y\_train}
  \item reentrena brevemente
  \item compara accuracy vs baseline de clase mayoritaria
  \item regla: \texttt{leakage\_detected = shuffled\_acc > baseline\_acc + 0.05}
\end{itemize}

Objetivo: comprobar que el rendimiento alto desaparece al romper la relación real X-y.

\subsubsection{5.3 Check C (split por día/CSV)}

Archivo: \verb|src/validate_checks.py|, función \verb|check_c_csv_split|

\begin{itemize}
  \item entrena en un conjunto de días/CSV
  \item evalúa en días/CSV diferentes
  \item reporta mismas métricas de clasificación + conteos
\end{itemize}

Objetivo: medir generalización fuera del split aleatorio.

\subsubsection{5.4 Leave-one-exact-CSV-out}

Archivo: \verb|src/validate_leave_one_csv_out.py|

\begin{itemize}
  \item cada fold deja un CSV oficial completo para test
  \item agrega: medias/desviaciones/min/max por métrica, matriz global agregada, métricas globales
        (accuracy, balanced accuracy, F1, FPR/FNR, reward)
\end{itemize}

Es la validación más estricta a nivel de separación por archivo real.

\subsection{6) Parámetros de recompensa (impacto en comportamiento)}

Valor por defecto común en entrenamiento/validaciones:

\begin{itemize}
  \item \texttt{tp=1.5}
  \item \texttt{fp=-1.5}
  \item \texttt{fn=-5.0}
  \item \texttt{omission=0.0}
\end{itemize}

Lectura para defensa:

\begin{itemize}
  \item se penaliza más el falso negativo que el falso positivo
  \item por diseño, el agente está sesgado a evitar dejar pasar ataques
\end{itemize}

\subsection{7) Qué puede preguntarte el tribunal y cómo enlazarlo}

\begin{enumerate}
  \item \textbf{``¿Cómo verificas que no hay leakage?''} --- Check~B (labels barajadas) + política
        explícita de drop de columnas de identificación.
  \item \textbf{``¿Cómo pruebas generalización realista?''} --- Check~C y
        leave-one-exact-CSV-out.
  \item \textbf{``¿Qué alternativa no-RL tienes?''} --- baseline Random Forest.
  \item \textbf{``¿Cómo soportas inferencia robusta en datos reales?''} --- pipeline v2 con
        clipping + diagnósticos de drift + artefactos reproducibles.
\end{enumerate}
